We validate the benchmark design with pilot studies on all five
suites---external, history, ICL, RAG, and tool---using unified
metrics (\S\ref{sec:metrics}) that enable direct cross-suite
comparison.

%% ----------------------------------------------------------
\subsection{Pilot setup}
\label{sec:pilot-setup}

\paragraph{External pilot.}
90~items from the external suite (15 per domain, 6~domains),
each rendered in 5~conditions (control, irrelevant\_low,
irrelevant\_high, plausible\_low, plausible\_high),
yielding 450~prompts.

\paragraph{History pilot.}
60~items from the history suite (10 per domain, 6~domains),
each rendered in 5~conditions using the two-stage protocol
(\S\ref{sec:suites}), yielding 300~prompts.
Non-control conditions require two inference stages per prompt:
Stage~1 produces the model's self-generated anchor, which is
then retained in conversation history for Stage~2.

\paragraph{RAG pilot.}
180~items from the RAG suite (30 per domain, 6~domains),
each rendered in 5~conditions using the three-document
corpus design (\S\ref{sec:suites}), yielding 900~prompts.
All conditions share the same evidence and question;
only the anchor-slot document varies.

\paragraph{ICL pilot.}
180~items from the ICL suite (30 per domain, 6~domains),
each rendered in 5~conditions using the metadata-priming
design (\S\ref{sec:suites}), yielding 900~prompts.
Three few-shot demonstrations precede the target item.
All conditions share the same target evidence, question,
demo evidence, and demo answers; only the demo header
metadata varies (case~IDs for irrelevant, prior estimates
for plausible, plain headers for control).

\paragraph{Tool pilot.}
180~items from the tool suite (30 per domain, 6~domains),
each rendered in 5~conditions using the structured tool-output
design (\S\ref{sec:suites}), yielding 900~prompts.
All conditions share the same evidence and
\texttt{get\_evidence\_summary} output; only the
\texttt{check\_external\_reference} result varies.
Prompts are rendered through each model's native chat template
with \texttt{role:~"tool"} messages.

\paragraph{Models.}
Four instruction-tuned models spanning two families and two
scale tiers:
Qwen2.5-7B-Instruct, Qwen2.5-3B-Instruct~\citep{qwen3},
Llama-3.1-8B-Instruct, and
Llama-3.2-3B-Instruct~\citep{llama3}.
All models are served on NVIDIA H100 GPUs with greedy decoding
(temperature $= 0.0$, max tokens $= 512$).

\paragraph{Parsing.}
Integer answers are extracted via a regex parser that locates
the last standalone integer in $[0,100]$.
When the regex fails, an LLM fallback
(Qwen2.5-1.5B-Instruct) extracts the final numeric answer.
Parse rates exceed 98\% on the external, history, ICL, and
RAG pilots for all models (the ICL pilot achieves 100\%
across all four models).
On the tool pilot, Qwen models maintain high parse rates
($\geq$98\%), while Llama models show reduced rates
(79--86\%), reflecting the additional challenge of generating
well-formatted answers after structured tool-calling
context.

%% ----------------------------------------------------------
\subsection{External pilot results}
\label{sec:pilot-results-external}

Table~\ref{tab:external_results} reports the unified metrics
for the external suite.

\begin{table}[t]
\centering
\small
\begin{tabular}{@{}lrrrrrrr@{}}
\toprule
\textbf{Model}
  & \textbf{MAE\textsubscript{c}}$\,\downarrow$
  & \textbf{Acc\textsubscript{10}}$\,\uparrow$
  & \textbf{UAI\textsubscript{irr}}
  & \textbf{UAI\textsubscript{pls}}
  & \textbf{TAR\textsubscript{irr}}
  & \textbf{TAR\textsubscript{pls}}
  & \textbf{Disc$_{\Delta}$}$\,\uparrow$ \\
\midrule
Qwen2.5-7B    & \textbf{7.18}  & \textbf{76.7\%} & 0.05  & 0.27 & 0.18 & 0.64 & 0.22 \\
Llama-3.1-8B  & 8.97  & 75.3\% & $-$0.22 & 0.43 & 0.32 & 0.74 & \textbf{0.66} \\
Llama-3.2-3B  & 10.54 & 66.7\% & 0.23  & 0.43 & 0.46 & 0.67 & 0.19 \\
Qwen2.5-3B    & 11.81 & 56.7\% & 0.05  & 0.30 & 0.27 & 0.54 & 0.25 \\
\bottomrule
\end{tabular}
\caption{External-pilot results (90~items, 5~conditions per
  item).  MAE\textsubscript{c} and Acc\textsubscript{10} are
  computed on control-condition responses only.
  UAI: Unified Anchor Influence; TAR: Toward-Anchor Rate;
  Disc$_{\Delta}$: UAI\textsubscript{pls} $-$
  UAI\textsubscript{irr}.
  $\varepsilon = 3$ for UAI filtering.}
\label{tab:external_results}
\end{table}

\paragraph{Accuracy and discrimination are different
  dimensions.}
Qwen2.5-7B (MAE~7.18, Acc~76.7\%) and Llama-3.1-8B
(MAE~8.97, Acc~75.3\%) have comparable task accuracy, yet
their anchoring profiles differ qualitatively.
Llama-3.1-8B achieves the highest discrimination gap
(Disc$_{\Delta} = 0.66$), driven by a mild contrast effect
under irrelevant anchoring (UAI\textsubscript{irr}~$= -0.22$)
combined with substantial plausible sensitivity
(UAI\textsubscript{pls}~$= 0.43$, TAR~$= 0.74$).
Qwen2.5-7B shows low anchoring overall
(UAI\textsubscript{irr}~$= 0.05$,
UAI\textsubscript{pls}~$= 0.27$) with moderate
discrimination (0.22).

\paragraph{TAR reveals directional effects that UAI magnitude
  can miss.}
Llama-3.2-3B shows substantial irrelevant influence
(UAI\textsubscript{irr}~$= 0.23$,
TAR\textsubscript{irr}~$= 0.46$) and plausible influence
(UAI\textsubscript{pls}~$= 0.43$,
TAR\textsubscript{pls}~$= 0.67$), yielding low
discrimination (Disc$_{\Delta} = 0.19$).
This indiscriminate anchoring pattern---high sensitivity
regardless of relevance---is invisible to accuracy metrics
alone: Llama-3.2-3B achieves better accuracy (66.7\%) than
Qwen2.5-3B (56.7\%).

%% ----------------------------------------------------------
\subsection{History pilot results}
\label{sec:pilot-results-history}

Table~\ref{tab:history_results} reports the unified metrics
for the history suite.
For this suite, the anchor value is the model's own Stage~1
answer rather than an experimenter-defined value
(\S\ref{sec:metrics}).

\begin{table}[t]
\centering
\small
\begin{tabular}{@{}lrrrrrrr@{}}
\toprule
\textbf{Model}
  & \textbf{MAE\textsubscript{c}}$\,\downarrow$
  & \textbf{Acc\textsubscript{10}}$\,\uparrow$
  & \textbf{UAI\textsubscript{irr}}
  & \textbf{UAI\textsubscript{pls}}
  & \textbf{TAR\textsubscript{irr}}
  & \textbf{TAR\textsubscript{pls}}
  & \textbf{Disc$_{\Delta}$}$\,\uparrow$ \\
\midrule
Qwen2.5-7B    & \textbf{6.92}  & \textbf{85.0\%} & $-$0.08 & 0.35 & 0.41 & 0.57 & 0.43 \\
Llama-3.1-8B  & 7.97  & 81.7\% & 0.15  & 0.17 & 0.37 & 0.45 & 0.02 \\
Llama-3.2-3B  & 9.88  & 73.3\% & 0.05  & 0.54 & 0.53 & 0.65 & \textbf{0.49} \\
Qwen2.5-3B    & 11.70 & 50.0\% & $-$0.17 & 0.25 & 0.42 & 0.58 & 0.41 \\
\bottomrule
\end{tabular}
\caption{History-pilot results (60~items, two-stage protocol).
  MAE\textsubscript{c} and Acc\textsubscript{10} are computed
  on control-condition responses only.
  Anchor values are model-generated (Stage~1 answers).
  ACR and RR are reported in
  Appendix~\ref{app:history-secondary}.}
\label{tab:history_results}
\end{table}

\paragraph{Self-generated anchoring patterns differ from
  external anchoring.}
Under the unified framework, the History suite reveals
anchoring patterns that diverge sharply from the External
suite.
Llama-3.1-8B, which achieves the highest External-suite
discrimination (Disc$_{\Delta} = 0.66$), shows near-zero
discrimination on History (Disc$_{\Delta} = 0.02$):
its UAI\textsubscript{irr} (0.15) and
UAI\textsubscript{pls} (0.17) are nearly equal, indicating
that it responds similarly to its own prior answers regardless
of whether they pertain to the same or a different case.

\paragraph{Smaller models show higher plausible sensitivity
  to self-generated anchors.}
Llama-3.2-3B achieves the highest UAI\textsubscript{pls}
across all models and suites (0.54,
TAR\textsubscript{pls}~$= 0.65$), indicating that its
Stage~2 answers shift substantially toward its biased
Stage~1 estimates---consistent with the insufficient
adjustment pattern predicted by \citet{epley2001putting}.
Combined with low irrelevant influence
(UAI\textsubscript{irr}~$= 0.05$), this yields the highest
History-suite discrimination (Disc$_{\Delta} = 0.49$).

%% ----------------------------------------------------------
\subsection{RAG pilot results}
\label{sec:pilot-results-rag}

Table~\ref{tab:rag_results} reports the unified metrics
for the RAG suite.

\begin{table}[t]
\centering
\small
\begin{tabular}{@{}lrrrrrrr@{}}
\toprule
\textbf{Model}
  & \textbf{MAE\textsubscript{c}}$\,\downarrow$
  & \textbf{Acc\textsubscript{10}}$\,\uparrow$
  & \textbf{UAI\textsubscript{irr}}
  & \textbf{UAI\textsubscript{pls}}
  & \textbf{TAR\textsubscript{irr}}
  & \textbf{TAR\textsubscript{pls}}
  & \textbf{Disc$_{\Delta}$}$\,\uparrow$ \\
\midrule
Llama-3.1-8B  & \textbf{6.58}  & \textbf{80.0\%} & 0.07 & 0.24 & 0.34 & 0.51 & 0.17 \\
Qwen2.5-3B    & 9.65  & 64.4\% & 0.00 & 0.10 & 0.18 & 0.38 & 0.10 \\
Llama-3.2-3B  & 9.97  & 67.2\% & 0.11 & 0.07 & 0.41 & 0.44 & $-$0.04 \\
Qwen2.5-7B    & 11.77 & 50.6\% & $-$0.01 & 0.22 & 0.13 & 0.51 & 0.23 \\
\bottomrule
\end{tabular}
\caption{RAG-pilot results (180~items, 5~conditions per
  item).  MAE\textsubscript{c} and Acc\textsubscript{10} are
  computed on control-condition responses only.}
\label{tab:rag_results}
\end{table}

\paragraph{Document context reverses the accuracy ranking.}
On the External suite, Qwen2.5-7B achieves the best task
accuracy (MAE~7.18); on the RAG suite, it is the worst
(MAE~11.77, Acc~50.6\%).
Conversely, Llama-3.1-8B improves from MAE~8.97 (External)
to MAE~6.58 (RAG), achieving the highest accuracy.
This reversal indicates that evidence aggregation in the
presence of retrieved documents is a distinct competence
from estimation in prompt-only settings.

\paragraph{RAG context attenuates anchoring overall but does
  not eliminate it.}
Across all models, UAI\textsubscript{pls} on RAG (range
0.07--0.24) is substantially lower than on External (range
0.27--0.43).
The document context---which provides domain-relevant
background alongside the anchor---appears to dilute anchor
influence.
However, Llama-3.2-3B shows negative discrimination on RAG
(Disc$_{\Delta} = -0.04$): it responds more to irrelevant
document metadata (UAI\textsubscript{irr}~$= 0.11$,
TAR~$= 0.41$) than to plausible survey results
(UAI\textsubscript{pls}~$= 0.07$,
TAR~$= 0.44$)---a concerning pattern for deployed RAG
systems.

%% ----------------------------------------------------------
\subsection{Tool pilot results}
\label{sec:pilot-results-tool}

Table~\ref{tab:tool_results} reports the unified metrics
for the tool suite, in which anchors appear within structured
JSON tool outputs rendered via the model's native chat template.

\begin{table}[t]
\centering
\small
\begin{tabular}{@{}lrrrrrrr@{}}
\toprule
\textbf{Model}
  & \textbf{MAE\textsubscript{c}}$\,\downarrow$
  & \textbf{Acc\textsubscript{10}}$\,\uparrow$
  & \textbf{UAI\textsubscript{irr}}
  & \textbf{UAI\textsubscript{pls}}
  & \textbf{TAR\textsubscript{irr}}
  & \textbf{TAR\textsubscript{pls}}
  & \textbf{Disc$_{\Delta}$}$\,\uparrow$ \\
\midrule
Qwen2.5-7B    & \textbf{0.54}  & \textbf{98.9\%} & 0.00  & 0.16 & 0.11 & 0.44 & 0.16 \\
Qwen2.5-3B    & 3.45  & 92.6\% & 0.06  & 0.09 & 0.24 & 0.29 & 0.04 \\
Llama-3.1-8B  & 15.04 & 56.3\% & 0.13  & 0.51 & 0.28 & 0.65 & \textbf{0.38} \\
Llama-3.2-3B  & 18.84 & 35.6\% & 0.22  & 0.39 & 0.38 & 0.55 & 0.17 \\
\bottomrule
\end{tabular}
\caption{Tool-pilot results (180~items, 5~conditions per
  item).  Anchors appear in structured \texttt{role:~"tool"}
  messages rendered through each model's chat template.
  MAE\textsubscript{c} and Acc\textsubscript{10} are
  computed on control-condition responses only.
  Parse rates: Qwen2.5-7B 100\%, Qwen2.5-3B 98.1\%,
  Llama-3.1-8B 79.3\%, Llama-3.2-3B 86.1\%.}
\label{tab:tool_results}
\end{table}

\paragraph{Structured tool context creates a sharp model-family
  divide.}
The tool suite reveals the starkest accuracy split of any suite:
Qwen models achieve near-perfect evidence aggregation
(MAE~0.54--3.45, Acc~93--99\%), while Llama models struggle
substantially (MAE~15--19, Acc~36--56\%).
This gap---absent on the external and history suites---suggests
that the ability to extract and aggregate numeric evidence from
structured \texttt{role:~"tool"} messages is a distinct
competence that varies across model families.
The lower parse rates for Llama models (79--86\% vs.\
98--100\% for Qwen) indicate that the tool-calling chat
template itself poses a formatting challenge.

\paragraph{Anchoring susceptibility and tool-context competence
  co-vary across model families.}
Llama-3.1-8B shows the highest Tool-suite anchoring
sensitivity
(UAI\textsubscript{pls}~$= 0.51$,
TAR\textsubscript{pls}~$= 0.65$) and the highest
discrimination gap (Disc$_{\Delta} = 0.38$).
In contrast, Qwen2.5-7B nearly ignores irrelevant tool metadata
(UAI\textsubscript{irr}~$= 0.00$) while remaining moderately
sensitive to plausible reference values
(UAI\textsubscript{pls}~$= 0.16$)---a desirable anchoring
profile for deployed agentic systems.
Qwen2.5-3B shows near-zero discrimination
(Disc$_{\Delta} = 0.04$), indicating that it responds
similarly to irrelevant and plausible tool outputs---not
because it ignores both, but because its low overall
anchoring sensitivity leaves little room for differentiation.
However, the large accuracy and parse-rate gap between
families means that anchoring susceptibility and
structured-context competence are confounded in this suite:
models that struggle to extract evidence from tool outputs
may also be more susceptible to anchor values embedded in
those outputs, and the present design cannot fully
disentangle these factors.

%% ----------------------------------------------------------
\subsection{ICL pilot results}
\label{sec:pilot-results-icl}

Table~\ref{tab:icl_results} reports the unified metrics
for the ICL suite, in which anchors appear only in the
metadata headers of few-shot demonstrations---case~IDs and
batch numbers for irrelevant conditions, prior estimates for
plausible conditions---while demo evidence and answers remain
identical across all conditions.

\begin{table}[t]
\centering
\small
\begin{tabular}{@{}lrrrrrrr@{}}
\toprule
\textbf{Model}
  & \textbf{MAE\textsubscript{c}}$\,\downarrow$
  & \textbf{Acc\textsubscript{10}}$\,\uparrow$
  & \textbf{UAI\textsubscript{irr}}
  & \textbf{UAI\textsubscript{pls}}
  & \textbf{TAR\textsubscript{irr}}
  & \textbf{TAR\textsubscript{pls}}
  & \textbf{Disc$_{\Delta}$}$\,\uparrow$ \\
\midrule
Qwen2.5-7B    & \textbf{5.79}  & \textbf{83.3\%} & $-$0.02 & $-$0.01 & 0.19 & 0.28 & 0.00 \\
Llama-3.1-8B  & 7.44  & 80.0\% & 0.14  & 0.19 & 0.48 & 0.52 & \textbf{0.05} \\
Qwen2.5-3B    & 7.56  & 74.4\% & $-$0.01 & $-$0.04 & 0.14 & 0.19 & $-$0.03 \\
Llama-3.2-3B  & 6.61  & 82.2\% & 0.04  & 0.09 & 0.40 & 0.48 & 0.05 \\
\bottomrule
\end{tabular}
\caption{ICL-pilot results (180~items, 5~conditions per
  item).  Anchors appear only in demo-header metadata;
  demo evidence and answers are identical across conditions.
  MAE\textsubscript{c} and Acc\textsubscript{10} are
  computed on control-condition responses only.
  Parse rate is 100\% for all models.}
\label{tab:icl_results}
\end{table}

\paragraph{Metadata priming produces the weakest anchoring
  of any suite.}
Across all five suites, the ICL metadata-priming design
yields the smallest UAI magnitudes.
Both Qwen models show near-zero or slightly negative UAI
(Qwen2.5-7B: UAI\textsubscript{irr}~$= -0.02$,
UAI\textsubscript{pls}~$= -0.01$; Qwen2.5-3B:
UAI\textsubscript{irr}~$= -0.01$,
UAI\textsubscript{pls}~$= -0.04$), indicating effective
immunity to incidental metadata numbers in demonstration
context.
Llama models show mild but detectable priming
(Llama-3.1-8B: UAI\textsubscript{pls}~$= 0.19$,
TAR\textsubscript{pls}~$= 0.52$), consistent with the
selective accessibility mechanism: metadata numbers
activate anchor-consistent retrieval that subtly influences
the estimate.

\paragraph{TAR confirms the absence of reliable toward-anchor
  shifts.}
TAR values reinforce the conclusion that metadata priming
produces at most fragile effects.
No model achieves TAR\textsubscript{irr} above the 0.50
chance level (Llama-3.1-8B: 0.48; Llama-3.2-3B: 0.40;
Qwen models: 0.14--0.19), indicating no systematic
toward-anchor shift under irrelevant metadata.
Only Llama-3.1-8B shows TAR\textsubscript{pls} marginally
above chance (0.52); the remaining models remain at or below
0.50.
The family-level pattern is nevertheless informative:
Qwen models show TAR well below chance on both relevance
levels, consistent with effective immunity to metadata
priming, while Llama models show near-chance TAR, indicating
at most borderline sensitivity that does not reliably exceed
chance.

%% ----------------------------------------------------------
\subsection{Cross-suite comparison}
\label{sec:cross-suite}

Table~\ref{tab:cross_suite} summarizes the key unified
metrics across all five suites.

\begin{table}[t]
\centering
\small
\begin{tabular}{@{}llrrrr@{}}
\toprule
\textbf{Suite} & \textbf{Model}
  & \textbf{MAE\textsubscript{c}}$\,\downarrow$
  & \textbf{UAI\textsubscript{irr}}
  & \textbf{UAI\textsubscript{pls}}
  & \textbf{Disc$_{\Delta}$}$\,\uparrow$ \\
\midrule
External & Qwen2.5-7B     & 7.18  & 0.05    & 0.27 & 0.22 \\
         & Llama-3.1-8B   & 8.97  & $-$0.22 & 0.43 & 0.66 \\
         & Llama-3.2-3B   & 10.54 & 0.23    & 0.43 & 0.19 \\
         & Qwen2.5-3B     & 11.81 & 0.05    & 0.30 & 0.25 \\
\midrule
History  & Qwen2.5-7B     & 6.92  & $-$0.08 & 0.35 & 0.43 \\
         & Llama-3.1-8B   & 7.97  & 0.15    & 0.17 & 0.02 \\
         & Llama-3.2-3B   & 9.88  & 0.05    & 0.54 & 0.49 \\
         & Qwen2.5-3B     & 11.70 & $-$0.17 & 0.25 & 0.41 \\
\midrule
ICL      & Qwen2.5-7B     & 5.79  & $-$0.02 & $-$0.01 & 0.00 \\
         & Llama-3.2-3B   & 6.61  & 0.04    & 0.09 & 0.05 \\
         & Llama-3.1-8B   & 7.44  & 0.14    & 0.19 & 0.05 \\
         & Qwen2.5-3B     & 7.56  & $-$0.01 & $-$0.04 & $-$0.03 \\
\midrule
RAG      & Llama-3.1-8B   & 6.58  & 0.07    & 0.24 & 0.17 \\
         & Qwen2.5-3B     & 9.65  & 0.00    & 0.10 & 0.10 \\
         & Llama-3.2-3B   & 9.97  & 0.11    & 0.07 & $-$0.04 \\
         & Qwen2.5-7B     & 11.77 & $-$0.01 & 0.22 & 0.23 \\
\midrule
Tool     & Qwen2.5-7B     & 0.54  & 0.00    & 0.16 & 0.16 \\
         & Qwen2.5-3B     & 3.45  & 0.06    & 0.09 & 0.04 \\
         & Llama-3.1-8B   & 15.04 & 0.13    & 0.51 & 0.38 \\
         & Llama-3.2-3B   & 18.84 & 0.22    & 0.39 & 0.17 \\
\bottomrule
\end{tabular}
\caption{Cross-suite comparison using unified metrics.
  Models are ordered by MAE\textsubscript{c} within each
  suite.  The same model can rank differently across suites
  on both accuracy and discrimination.}
\label{tab:cross_suite}
\end{table}

\paragraph{Anchoring susceptibility is paradigm-dependent.}
The five pilots collectively demonstrate that anchoring is
not a unitary property of a model but depends on the delivery
mechanism.
Llama-3.1-8B exhibits the highest External-suite
discrimination (Disc$_{\Delta} = 0.66$) but near-zero
History-suite discrimination (0.02), moderate RAG and
Tool-suite discrimination (0.17 and 0.38), and small
ICL-suite discrimination (0.05).
Llama-3.2-3B achieves the highest History-suite
discrimination (0.49) but the lowest External-suite
discrimination (0.19) and negative RAG-suite discrimination
($-$0.04).
No model is consistently best or worst across paradigms.

\paragraph{The accuracy ranking reverses across suites.}
The model ranking by task accuracy is consistent across
External and History suites
(Qwen2.5-7B $>$ Llama-3.1-8B $>$ Llama-3.2-3B $>$
Qwen2.5-3B) but reverses on the RAG suite, where
Llama-3.1-8B leads and Qwen2.5-7B trails.
The Tool suite introduces a third ranking pattern:
Qwen models dominate (MAE~0.54--3.45) while Llama models
lag substantially (MAE~15--19).
The ICL suite introduces yet another ranking:
Qwen2.5-7B leads (MAE~5.79) while Llama-3.2-3B outperforms
both Llama-3.1-8B and Qwen2.5-3B, suggesting that evidence
aggregation competence is context-format-dependent.

\paragraph{Plausible anchoring is consistently stronger than
  irrelevant anchoring.}
Across all 20 model--suite combinations,
UAI\textsubscript{pls} exceeds UAI\textsubscript{irr} in
18 of 20 cases (Disc$_{\Delta} > 0$), with two exceptions:
Llama-3.2-3B on RAG ($-$0.04) and Qwen2.5-3B on ICL
($-$0.03).
This overall pattern is encouraging---models generally
\emph{do} distinguish relevance levels---but the
magnitude and consistency of discrimination vary
substantially across paradigms, confirming the value of
multi-suite evaluation.

\paragraph{Metadata priming is the subtlest anchoring
  pathway.}
The ICL suite yields the smallest anchoring effects of any
suite: UAI\textsubscript{pls} ranges from $-0.04$ to $0.19$,
compared to $0.07$--$0.54$ across the other four suites.
This is consistent with the selective accessibility theory:
incidental metadata numbers in demonstration headers are a
weaker prime than explicit anchor sentences (External),
self-generated estimates (History), document content (RAG),
or structured tool outputs (Tool).
Qwen models are effectively immune to metadata priming
(UAI near zero, TAR well below chance), while Llama models
show near-chance TAR values (0.40--0.48 for irrelevant),
indicating at most borderline directional sensitivity that
does not reliably exceed the 0.50 chance level.

\paragraph{Structured tool outputs amplify model-family
  differences.}
The Tool suite's accuracy gap between model families
(Qwen MAE~0.5--3.5 vs.\ Llama MAE~15--19) far exceeds the
spread observed on any other suite, where all four models
cluster within a 5--8-point MAE range.
At the same time, Tool-suite \emph{anchoring} mirrors the
External suite more closely than it does the RAG or ICL
suites:
Llama-3.1-8B achieves high discrimination on both External
(0.66) and Tool (0.38), while Qwen2.5-3B shows low
discrimination on both (0.25 and 0.04).
This partial correlation suggests that external and
tool-output anchoring may share a susceptibility pathway,
while document-based and metadata-priming anchoring engage
different processing modes---though the confound between
anchoring susceptibility and tool-context competence
(\S\ref{sec:pilot-results-tool}) warrants caution in
interpreting the Tool-suite patterns.
